\documentclass[11pt]{article}
\usepackage{amsmath,amssymb,amsthm}
\usepackage{geometry}
\geometry{margin=1in}
\newtheorem{theorem}{Theorem}
\newtheorem{definition}{Definition}

\title{Refutation of Conjecture 00000004413\\
\large Cycle counts in bounded-degree graph limits}
\author{orionsheep (TLMC refutation submission)}
\date{October 3, 2026}

\begin{document}
\maketitle

\section{The conjecture}

\begin{definition}[graphing / bounded-degree graph limit]
A (deterministic) graphing on a finite probability space $V$ is a symmetric,
irreflexive edge relation $E$ on $V$.  It satisfies \emph{degree bound} $d$
when every vertex has at most $d$ neighbours.  Its \emph{cycle measure} of
length $l$ is the number of distinct undirected $l$-cycles divided by $|V|$
(the uniform measure on $V$).
\end{definition}

\noindent\textbf{Conjecture 00000004413.} \emph{For every graphing of degree
bound $d$ and every $l$, the $l$-cycle measure is at most
$\dfrac{(d-1)^l}{2l}$, and this bound is attained and optimal for the
$d$-regular tree limit.}

\section{The conjecture is false}

\subsection{Counterexample to the upper bound}

Take $G = K_3$, the triangle, on the probability space $V = \{0,1,2\}$
($x\sim y \iff x\neq y$).

\begin{itemize}
  \item Degree bound: every vertex has degree $2$, so $d = 2$.
  \item The number of undirected $3$-cycles is $1$, so the $3$-cycle measure
        is $1/3$.
  \item The claimed bound: $\dfrac{(d-1)^l}{2l} = \dfrac{1^3}{6} = \dfrac16$.
\end{itemize}

\noindent But $\dfrac13 > \dfrac16$: the alleged upper bound fails on the
simplest bounded-degree object.  Already the raw cycle \emph{count} $1$
exceeds $1/6$.

\subsection{The optimality/attainment claim is also false}

A $d$-regular \emph{tree} contains no cycles at all, so its $l$-cycle
measure is $0$, whereas the claimed bound $\frac{(d-1)^l}{2l}$ is strictly
positive for every $d \ge 2$ and $l \ge 1$.  Hence the bound is never
``attained'' by the $d$-regular tree: the second half of the conjecture is
inconsistent with the first.

\subsection{What $(d-1)^l/(2l)$ really is}

The expression $(d-1)^l/(2l)$ is the \emph{mean of the Poisson limit} for
the number of $l$-cycles in a uniformly random $d$-regular graph on $n$
vertices (Bollob\'as--Wormald): the cycle count of a random $d$-regular
graph converges in distribution to $\mathrm{Poisson}\bigl((d-1)^l/(2l)\bigr)$.
It is therefore a \emph{typical/expected} value for locally-tree-like limits,
not a universal upper bound.

\section{Formalization}

The Lean~4/Mathlib development (\texttt{lean4/Main.lean}) defines

\begin{itemize}
  \item \texttt{Graphing V}: symmetric irreflexive decidable relation on a
        finite vertex type $V$ (uniform measure);
  \item \texttt{Graphing.deg}, \texttt{Graphing.DegreeBound},
        \texttt{Graphing.cycleWalks}, \texttt{Graphing.edgeSetOf}
        (a \texttt{Finset (Sym2 V)}, so a cycle and all its
        rotations/reversal give the same edge set),
        \texttt{Graphing.cycleCount} (normalized by $|V|$), and
        \texttt{Graphing.claimedBound} $= (d-1)^l/(2l)$;
  \item \texttt{triangle}: the $K_3$ graphing on \texttt{Fin 3}.
\end{itemize}

\noindent Main theorems:

\begin{itemize}
  \item \texttt{conjecture\_00000004413\_false}: the universal bound is
        \emph{false} (proved from
        \texttt{triangle\_degreeBound}, \texttt{triangle\_cycleCount\_ge}
        $\Rightarrow$ \texttt{cycleCount}~$\ge 1/3$, and
        \texttt{claimedBound\_2\_3} $= 1/6$);
  \item \texttt{counterexample\_exists}: explicit counterexample;
  \item \texttt{tree\_cannot\_attain}: $\texttt{claimedBound} > 0$ for
        $d \ge 2$, $l \ge 1$, while a tree has cycle measure $0$, so the
        bound cannot be attained by the $d$-regular tree.
\end{itemize}

\noindent The axiom audit (\texttt{lean4/Check.lean}, \texttt{\#print axioms})
reports dependence only on Lean's three standard foundational axioms
\texttt{propext}, \texttt{Classical.choice}, \texttt{Quot.sound} --- no
\texttt{sorry}, no \texttt{native\_decide}.

\section{Reproduction}

\begin{itemize}
  \item \texttt{reproduce.py}: brute-force enumeration of undirected
        $3$-cycles in $K_3$ ($=1$), comparison with $1^3/6 = 1/6$, and a
        check that trees have cycle count $0 < (d-1)^l/(2l)$.
  \item \texttt{lean4/}: \texttt{lake build} exits $0$ on Lean
        \texttt{v4.33.1}/Mathlib \texttt{v4.33.1}.
\end{itemize}

\end{document}
